\subsection{Limitations}
\label{sec:limitations}

Several limitations bound how far our ranking can be trusted. First, the wedge is
anchored to the affected \emph{plant} of each redispatch event, which is the
actuator the TSO uses, not the location of the binding constraint; the PTDF kernel
mitigates but cannot relocate a signal that is already mislocated, so scores reward
proximity to dispatch activity rather than to congestion itself. Second, our data
covers only the 220/380~kV transmission system, so a battery that relieves
distribution-level congestion can score poorly despite being well placed, making
the ranking valid only within the transmission frame. Third, roughly 30\% of
redispatch volume never enters the scoring---mostly countertrading, which by
construction identifies no congested location, and secondarily plants that cannot
be matched to coordinates---so the effective price is built on a partial view of
the system. Crucially, this missing share is unlikely to be random: matchability
correlates with plant type, region, and vintage, and countertrading concentrates on
particular corridors, so the missing events are plausibly clustered in space and the
recovered hotspot map may differ from the true one even at the transmission level we
do observe. Finally, batteries co-located with PV or wind can be good for the grid
almost by construction, absorbing local generation that would otherwise be curtailed,
yet our pure arbitrage score does not credit this avoided curtailment and so
under-ranks them. Taken together, the scores are best read as a relative,
transmission-level indication of placement against the redispatch we can observe and
localise, not an absolute, node-level verdict on any individual site.
